\subsubsection{Handle cohomology and marking coordinates}
\label{subsec:handle-cohomology-markings}

\begin{definition}[Cohomological marking and handle cocycle]
\label{def:handle-cocycle}
Let \((\mathcal H,\beta)\) be a marked Kirby datum.  Dualize the handle
chain complex of \cref{def:derived-handle-chain-complex}:
\[
  C^k(\mathcal H):=\operatorname{Hom}(C_k(\mathcal H),\mathbb Z),
  \qquad
  \delta^k(\omega):=\omega\circ\partial_{k+1},
\]
and write
\[
  H^2(\mathcal H):=
  \ker(\delta^2)/\operatorname{im}(\delta^1).
\]
The cellular comparison isomorphism
\[
  \kappa_{\mathcal H}\colon
  H^2(\mathcal H)\xrightarrow{\ \cong\ }
  H^2(|\mathcal H|;\mathbb Z)
\]
and Poincar\'e duality determine the class
\[
  \xi_{\mathcal H,\beta}
  :=\kappa_{\mathcal H}^{-1}
    \bigl(\operatorname{PD}_{|\mathcal H|}(\beta)\bigr)
  \in H^2(\mathcal H).
\]
A \emph{handle cocycle} for \((\mathcal H,\beta)\) is a cocycle
\(\omega\in C^2(\mathcal H)\) representing \(\xi_{\mathcal H,\beta}\).
If \(e_j\) is the oriented (2)-handle generator corresponding to the
\(j\)th framed component, set
\[
  \omega_j:=\omega(e_j)\in\mathbb Z.
\]
\end{definition}

\begin{proposition}[Cocycle coordinates under Kirby--Cerf moves]
\label{prop:marking-coordinate-transport}
Let \(m\colon(\mathcal H,\beta)\to
(\mathcal H',|m|_*\beta)\) be a marked elementary move, and let \(\omega\)
be a handle cocycle for \((\mathcal H,\beta)\).  The target class has a
representative \(\omega'\) with the following coordinates.
\begin{enumerate}[label=(\roman*),ref=\roman*]
  \item For framed isotopy, a \(1/1\) slide, a \(2/1\) passage, a
  \(3/3\) slide, a \(0\)--\(1\) birth or death, or a \(3\)--\(4\) birth or
  death, the \(2\)-handle generators are canonically identified and
  \(\omega'=\omega\).

  \item For a \(2/2\) slide, let \(e_x\) be the moving generator,
  \(e_y\) the stationary generator, and let
  \(\varepsilon\in\{\pm1\}\) be fixed by the oriented handle-slide
  convention.  Then
  \begin{equation}
    \omega_x'=\omega_x+\varepsilon\omega_y,
    \qquad
    \omega_y'=\omega_y,
    \label{eq:two-two-marking-transport}
  \end{equation}
  while every uninvolved coordinate is unchanged.

  \item A \(1\)--\(2\) birth admits \(\omega'=\omega\oplus0\); every
  representative of the target class is cohomologous to one whose new
  \(2\)-handle coordinate is zero.

  \item For a \(2\)--\(3\) birth, the cocycle equation forces the new
  \(2\)-handle coordinate to be zero, so again
  \(\omega'=\omega\oplus0\).
\end{enumerate}
The death formulas are the inverses of the corresponding birth formulas.
\end{proposition}

\begin{proof}
Poincar\'e duality is natural, so the class attached to \(|m|_*\beta\) is
the cellular transport of the class attached to \(\beta\).  The standard
cellular map of each local move gives the formulas above: a \(2/2\) slide
sends the new moving generator to \(e_x+\varepsilon e_y\), while an
adjacent-index birth adds the corresponding acyclic two-term complex.  See
\cite[Sections~4.1--4.2]{DurstGeigesKegel2019}.
\end{proof}

\begin{lemma}[Cochain normal forms for cancelling pairs]
\label{lem:cancellation-cochain-normal-forms}
Let \(\mathcal H\to\mathcal H^+\) be a birth in the isolated standard local
form of \cref{prop:local-movie-normal-forms}.  For a $1$--$2$ birth,
extension by zero on the new primal $2$-handle sends a source $2$-cocycle
to a target $2$-cocycle representing the transported class.  Conversely,
every target $2$-cocycle is cohomologous to one whose new $2$-handle
coordinate is zero.  The same statement applies to an isolated $1$--$2$
death after the slide reduction in \cref{item:cancellation-normal-form}.
For a standard $2$--$3$ birth or death, every cocycle on the side containing
the pair has zero coordinate on the added $2$-handle.
\end{lemma}

\begin{proof}
After the other incidence entries have been cleared by the $2/2$ slides in
\cref{item:cancellation-normal-form}, the handle cochain complex of the
isolated cancelling $1$--$2$ pair contains the direct summand
\[
  \mathbb Z\xrightarrow{\,\pm1\,}\mathbb Z
\]
in degrees $1$ and $2$.  Extension by zero is the inclusion of the
complementary cochain complex.  Subtracting the coboundary of an appropriate
multiple of the new degree-$1$ basis cochain makes the new degree-$2$
coordinate zero without changing its cohomology class.

For a cancelling $2$--$3$ pair the corresponding direct summand lies in
degrees $2$ and $3$, again with differential $\pm1$.  The cocycle equation
in degree $3$ therefore forces the new degree-$2$ coordinate to vanish.
\end{proof}
